\begin{table}[H]
\centering
\caption{Standardized Effect Sizes for Main Outcomes}
\label{tab:sde}
\begin{threeparttable}
\begin{tabular}{llccccc}
\toprule
Outcome & $\hat{\beta}$ & SD($X$) & SD($Y$) & SDE & SE(SDE) & Classification \\
\midrule
\multicolumn{7}{l}{\textit{Panel A: Pooled}} \\
HHI & 0.0120 & 0.213 & 0.246 & 0.010 & 0.105 & Small positive \\
Top-country share & 0.0087 & 0.213 & 0.227 & 0.008 & 0.093 & Small positive \\
N sources ($>$1\%) & -1.546 & 0.213 & 3.744 & -0.088 & 0.091 & Moderate negative \\
\midrule
\multicolumn{7}{l}{\textit{Panel B: Heterogeneous (China-dependent vs.~not)}} \\
HHI (China-dep.) & -0.0286 & 0.342 & 0.270 & -0.036 & 0.296 & Small negative \\
HHI (Non-China) & 0.0109 & 0.164 & 0.228 & 0.008 & 0.197 & Small positive \\
\bottomrule
\end{tabular}
\begin{tablenotes}[flushleft]
\small
\item \textit{Notes:} \textbf{Country:} European Union (EU-27). \textbf{Research question:} Whether the EU Critical Raw Materials Act (2024/1252), which imposes a 65\% single-country import concentration ceiling for strategic minerals, caused more-concentrated minerals to diversify their import sources. \textbf{Policy mechanism:} The CRMA mandates that no single third country may supply more than 65\% of EU consumption for 17 strategic minerals by 2030, creating regulatory pressure on firms to find alternative suppliers for heavily concentrated materials. \textbf{Outcome definition:} Herfindahl--Hirschman Index (HHI) of import source concentration, calculated from bilateral trade values as $\sum_i s_i^2$ where $s_i$ is a partner country's share. \textbf{Treatment:} Continuous --- pre-CRMA (2022) HHI interacted with post-proposal indicator (2023--2024). \textbf{Data:} UN Comtrade bilateral trade flows (HS4--HS6), 2018--2024 annual, EU-27 imports from all partners. \textbf{Method:} Continuous-treatment DiD with mineral and year fixed effects, mineral-clustered standard errors. \textbf{Sample:} 20 mineral commodities (16 CRMA strategic + 4 controls) across 7 years. SDE $= \hat{\beta} \times \text{SD}(X) / \text{SD}(Y)$ where SD($X$) is the standard deviation of the treatment variable and SD($Y$) is the pre-treatment standard deviation of HHI. Classification refers to magnitude, not statistical significance: Large ($|$SDE$| > 0.15$), Moderate ($0.05$--$0.15$), Small ($0.005$--$0.05$), Null ($< 0.005$).
\end{tablenotes}
\end{threeparttable}
\end{table}
